We presented a continuous hidden Markov model framework that learns Gaussian emission parameters and transition probabilities directly from observed financial time series via the Baum-Welch algorithm, augmented with a Poisson jump-duration mechanism for volatility clustering.
The key finding is that continuous Gaussian emissions achieve high distributional fidelity (87--95\% KS pass rates across $K \in \{3, 6, 9, 11, 13\}$ states) while eliminating the quantization error inherent in discrete state partitioning, though at the cost of systematic kurtosis underestimation relative to the Student-$t$ emissions of the discrete framework \citep{alswaidan2026hybrid}.

The state resolution analysis demonstrated that distributional fidelity improves monotonically with $K$, while the jump mechanism's temporal contribution (ACF-MAE) remains stable across all resolutions.
Cross-asset generalization to NVDA, JNJ, and JPM confirmed that the Baum-Welch training captures diverse equity volatility profiles without modification.
Out-of-sample evaluation on held-out 2025 data showed robust generalization, with KS pass rates comparable to in-sample performance.

Extension of the framework to VIX log returns---an asset class with fundamentally different statistical properties (positive skewness, mean reversion, extreme kurtosis)---demonstrated that the CHMM generalizes beyond equity returns to volatility dynamics.
The VIX analysis achieved KS pass rates above 99\% at higher state resolutions ($K \geq 9$) and provided a clearer picture of the jump mechanism's role: at small $K$, the base model captures volatility clustering through naturally broad emission states, but as $K$ increases, finer state granularity reduces tail-state residence times, making the Poisson jump extension essential for preserving realistic tail persistence.
This finding reframes the jump mechanism as a structural complement to the state resolution, not a correction for model misspecification.

Several directions for future work follow from these results.
First, replacing Gaussian emissions with Student-$t$ or generalized hyperbolic distributions within the Baum-Welch loop would address the kurtosis gap while retaining the continuous estimation framework.
Second, the regime-dependent emission variances learned from VIX data provide a natural, data-driven parameterization for the mean-reversion target $\theta(S_t)$ in a modified Heston stochastic volatility model \citep{heston1993closed}, enabling closed-form implied volatility generation as an output of the HMM regime structure rather than requiring market-observed IV surfaces as input.
This Heston integration would produce a fully structural pipeline---from observed prices through HMM regime identification to implied volatility and ultimately option pricing via binomial trees---without circular dependence on market option data.
Third, extending the framework to multi-asset portfolio generation via copula-based dependence models \citep{alswaidan2026hybrid} would enable full cross-sectional simulation while preserving each asset's CHMM marginal.
Finally, time-varying transition matrices estimated via rolling windows or Bayesian online learning would allow the model to adapt to structural regime shifts.

\paragraph{Data Availability.}
The simulation scripts, training and test data, and all code to reproduce the results in this paper are available under an MIT license from the paper repository: \url{https://github.com/altashly1/CHMM-paper}.
The VIX analysis code and results are available at: \url{https://github.com/cadejin/HMM-Vol}.
The core continuous HMM implementation is available in the development repository: \url{https://github.com/altashly1/HMM-withJumps-WIP}.

\paragraph{Conflict of Interest.}
The authors declare no competing interests.

\paragraph{Author Contributions.}
J.V.\ directed the study.
A.A.\ developed the continuous model, implemented the Baum-Welch algorithm, conducted the SPY analysis, and generated the SPY figures.
C.J.\ conducted the VIX analysis and generated the VIX figures.
All authors edited and reviewed the final manuscript.
